% !PRINTED 488
\weberpage{488}
From (10) there follows still further:
\weq{12}{f_2(\omega)f_2\!\left(\dfrac{\omega}{n}\right)=\sqrt2\,e^{-\frac{r\pi i}{12}},}
\weq{13}{f(\omega)f\!\left(\dfrac{\omega}{n}\right)f_1(\omega)f_1\!\left(\dfrac{\omega}{n}\right)=\sqrt2\,e^{\frac{r\pi i}{12}}.}
If we assume $n=23$ and $r=0$, then equation \S\,75, (8) gives, with the use
of (12):
\weq{14}{f(\omega)f\!\left(\dfrac{\omega}{n}\right)-f_1(\omega)f_1\!\left(\dfrac{\omega}{n}\right)=2+\sqrt2,}
from which, by squaring twice and using (11) and (13), it follows that:
\[
x^3+\frac{1}{x^3}=36+26\,\sqrt2;
\]
from this one derives the simpler equation:
\weq{15}{x+\dfrac{1}{x}=3+\sqrt2,\quad \sqrt2\,x=f_1^2(\sqrt{-46}).}
3. Finally we make one further application of the transformation equations
of \S\,76 for a composite degree of transformation, to the determinant
$-39$.

If we set
\[
u=f(\sqrt{-39}),\quad v=f\!\left(\sqrt{\dfrac{-13}{3}}\right),\quad \frac{u^2+v^2}{uv}=z,
\]
then in equation (22), \S\,76, which relates to $n=39$, we have to set:
\[
A=\frac{u^2}{v^2}-2\,\frac{v}{u},\quad B=\frac{v^2}{u^2}-2\,\frac{u}{v},
\]
and the equation mentioned gives:
\[
z^3+z^2-5\,z-6=0,
\]
from which, after casting off the factor $z+2$, the following arises:
\weq{16}{z^2-z-3=0,\quad z=\dfrac{1+\sqrt{13}}{2}.}
Between $u$ and $v$, however, there subsists on the other hand a
transformation equation of the third order (\S\,73):
\[
\frac{u^6v^6-8}{u^3v^3}=\frac{u^{12}+v^{12}}{u^6v^6}=z^6-6\,z^4+9\,z^2-2,
\]
which by means of (16) can be brought to the form $z^4-2=\dfrac{27+7\sqrt{13}}{2}$.
